\subsection{Overview}

The simulation, design and performance evaluation of agricultural subsurface (tile) drainage rest on several converging lines of evidence: site-scale field validations of water-table depth, drain discharge, salinity and solute fluxes; benchmark comparisons of analytical drain-spacing solutions against experiments, numerical codes and one another; pipe- and conduit-scale hydraulic experiments; watershed- and regional-scale studies of how drainage is conceptualised within distributed models; formal sensitivity, calibration and uncertainty analyses; and a large body of tracer, hydrograph-separation and monitoring work that quantifies how drainage reshapes streamflow. Across this literature a robust set of consolidated findings has emerged. The field-scale water-management model \textsc{DRAINMOD} is by far the most extensively tested and is repeatedly judged a reliable design and evaluation tool for poorly and artificially drained lands \parencite{Singh2019b,Gowda2012}. Model skill is strongly scale-dependent in time: monthly and annual drainage volumes and cumulative water balances are reproduced with high accuracy, whereas daily and sub-daily peak flows, low dry-season discharge and snowmelt events are the persistent weak points \parencite{Gowda2012,Dayyani2009,Morrison2014,Varvaris2018}. Treating drainage as purely saturated flow can over-predict discharge by up to an order of magnitude, so capillary and unsaturated-zone storage cannot be neglected \parencite{Xin2016b}. Saturated hydraulic conductivity and the evapotranspiration algorithm dominate model sensitivity and uncertainty \parencite{NabiSichani2012,Malakshahi2020,Wang2006}. Unsteady-state spacing equations yield narrower, and evapotranspiration-aware solutions wider, more economical spacings than classical steady-state theory \parencite{OOSTERBAAN1988,NIKAM1992}. At pipe scale, entrance resistance and effective radius are now quantified for a wide range of perforation and envelope geometries \parencite{Ghane2022,Ghane2022b}. At the catchment scale, tile drainage supplies a major share of streamflow (frequently 40--65\%), raises low flows and baseflow, damps intermediate flows, has modest effect on the largest floods, and homogenises hydrologic response, while how drains are represented can shift the simulated drain contribution to streamflow by a factor of two or more \parencite{King2014,Schilling2008,Sloan2017,BolandBrien2014,Karan2021,Boico2022}. Data-driven, remote-sensing and geostatistical methods now complement process models for prediction, mapping and design. This section synthesises the quantified performance, mechanisms and design implications these studies have established.

\subsection{Field-Scale Water-Management Models}

\subsubsection{DRAINMOD as the Benchmark Field-Scale Model}

\textsc{DRAINMOD} is one of the first widely adopted and most extensively used models for assessing drainage performance in tile-drained fields, and its process-based structure with independently measurable inputs allows good description of drainage hydrology across a wide range of climates when soil hydraulic properties are adequately characterised \parencite{Singh2019b,Gowda2012}. Across humid and temperate sites it reproduces water-table depth and drainage discharge closely, with accuracy improving markedly from the daily to the monthly and annual scale. In North Carolina, water-table-depth Nash--Sutcliffe efficiencies (EF) were 0.68--0.72, daily drainage EF 0.49--0.73, monthly drainage EF 0.77--0.87, and annual drainage volumes lay within 3~cm (5\%) of observations, with a two-year cumulative difference of 1~cm (1\%) \parencite{Gowda2012}. In Iowa, cumulative 1990--2003 drainage was within 2--10\% of measured values with a slight tendency to over-predict in low-drainage years \parencite{Singh2006b}, and two-year eastern-Canada evaluations reported simulated annual drainage about 30~mm below observed while capturing event timing and magnitude \parencite{Dayyani2009}. A recurring tendency is to underestimate daily peak flows even where general trends are captured \parencite{Morrison2014,Singh1996}. Uncertainty analysis in Indiana confirmed the scale dependence: 90\% confidence-interval widths grew from 0--0.6~cm~day$^{-1}$ (daily) to 7.6--12.4~cm~year$^{-1}$ (yearly), yet annual predictions fell well within the bounds \parencite{Wang2006}.

\textsc{DRAINMOD} has been validated well beyond its original humid setting. In cold eastern Canada it accurately simulated subsurface-drainage and water-table timing and magnitude including winter and spring snowmelt ($R^2$ up to 0.93 for water-table depth), with calibration improving monthly performance but surface runoff and daily peaks predicted less well, and better performance in uniform than in cracking soils \parencite{Dayyani2009,Morrison2014}. Under semi-arid Indian conditions, midspan water-table heights agreed to average absolute deviations of 14--24~cm, though the model responded too fast to rainfall and over-estimated depths in dry and under-estimated them in wet periods, yet was judged satisfactory for drainage design \parencite{GUPTA1993}. On Illinois claypan soils weekly cumulative SEW30 agreed at $R^2=0.91$, and subsurface drainage affected trafficability and SEW30 far more than surface drainage \parencite{MOSTAGHIMI1989}; further acceptable-to-good calibrations ($R^2=0.89$--0.93, NSE $=0.84$--0.89) have been reported for groundwater-table prediction \parencite{Abduljaleel2023,Awad2022b}.

The literature also delineates failure modes. At an Israeli site water-table differences reached 0.3--1.7~m in four of five replicates, attributed to a river boundary, permeable sandy layers below drain depth, and deep-plus-lateral seepage exceeding 50\% of annual irrigation plus rainfall \parencite{Sinai2006}. A structural drawback is the inability to represent time-variable surface depression storage, producing an unrealistic water fate and a marked observed--predicted gap during a mid-summer paddy drainage period \parencite{Abduljaleel2023}. The model over-predicts drainage in low-drainage years \parencite{Singh2006b} and, being built around parallel-drain geometry, requires plot-by-plot calibration in irregular landscapes, where equivalent drain-spacing parameters ranged from about 1700 to 4000~cm and organic layers strongly influenced water tables \parencite{He2002}.

Because site-specific hydraulic data are often lacking, pedotransfer estimation combined with automated calibration has been shown to drive \textsc{DRAINMOD} reliably. ROSETTA-derived properties with the PEST estimator produced calibrated Iowa applications \parencite{Singh2006b}, and three levels of soil hydraulic parameter sets all gave acceptable annual and monthly drainage, although calibrated parameters best in calibration performed worst in validation \parencite{Qi2015}. In South Africa, van Genuchten fits ($R^2>0.90$) supported validation NSE of 0.86 (water table) and 0.81 (discharge) \parencite{Malota2024,Malota2015}. Enhancements target the daily-peak weakness: incorporating macropore flow captured about 75\% of 60 peak events and predicted cumulative drainage within 8\%, though surface runoff was strongly over-predicted \parencite{Askar2020}. A neural-network surrogate (\textsc{DRAINMOD-ANN}) cloned daily water-table predictions across layouts without altering source code, GRNN outperforming feed-forward training \parencite{Awad2022,Awad2022b}. Sensitivity to the evapotranspiration algorithm is pronounced: replacing the default Thornthwaite method with Rohwer or pan-based estimates substantially improved water-table prediction in Iranian paddies, and the best ET methods were not necessarily best for water-table simulation \parencite{Malakshahi2020}; moving ET$_0$ from a daily to an hourly step shifted the water fate, raising runoff 4.8\% and lowering drainage 3.1\% \parencite{Awad2022b}.

\subsubsection{Variably Saturated, Comparative and Specialised Field Models}

Richards-equation codes complement \textsc{DRAINMOD}. HYDRUS (2D/3D) consistently predicts drainage fluxes and water-table depths under contrasting systems (EF 0.91--1; RMSE 0.2--6.4~cm) in Iranian paddies and is recommended for identifying optimal layouts \parencite{Karandish2018,DarziNaftchali2018}. It reproduced Canadian soil-water content (67--76\% of variance) and clarified mechanics, showing that in compacted histosols matric-potential drawdown was 58\% faster below than above a compact layer, so raising that layer's conductivity (by up to $50\times$) rather than narrowing spacing is the effective remedy \parencite{Mante2018,Guedessou2021}; it also simulated higher drainage than the Hooghoudt and Murashima--Ogino equations for paddies and highlighted the control of spacing, depth and cracks \parencite{Ebrahimian2015}. Related variably saturated codes include the CSUID 3D model, which reproduced 1D infiltration, 2D water-table and drainage benchmarks ($R^2=0.96$--0.99) while propagating parameter uncertainty to root-zone salinity and yield \parencite{Alzraiee2013b,Alzraiee2013}, and Richards/Darcy simulations of sheet-pipe systems showing strong textural control of drawdown \parencite{Saptomo2021,Li2021}.

Head-to-head comparisons generally find several field models perform comparably. In the Netherlands, WAVE, SWAP and \textsc{DRAINMOD} predicted lateral drainage flux equally reliably, \textsc{DRAINMOD} giving the most accurate discharge \parencite{ElSadek2001}. \textsc{DRAINMOD} and SWACROP predicted midspan water tables and outflow close to measurements without calibration, \textsc{DRAINMOD} slightly better, though both over-predicted outflow after dry spells \parencite{Prasher1996}. SWAP reproduced water-table behaviour and cumulative discharge in India (0.17--0.18~m; 7.8--39.4~mm) and was endorsed as a simple performance tool \parencite{Verma2014}. \textsc{WaSim} predicted water-table depth well ($R^2=0.86$) but discharge less confidently ($R^2=0.57$) under Rosetta-estimated conductivity \parencite{Malota2022}, yet was judged a workable design and reclamation tool elsewhere \parencite{Hirekhan2007,Malota2024}. Among distributed models, \textsc{DRAINMOD} tended to over-predict and CATHY to under-predict tile flow, both poor at low flow but good in wet periods, with the layer below the drains most influential \parencite{Muma2017}. For salinity, a saturated--unsaturated ``SD'' approach outperformed \textsc{DRAINMOD-S}, which over-estimated short-term and under-estimated long-term salinity \parencite{Nozari2021}, although \textsc{DRAINMOD-S} itself calibrated satisfactorily in coastal China \parencite{Feng2018}.

Broader agricultural models have been adapted with generally acceptable results. \textsc{RZWQM}/\textsc{RZWQM2} simulated weekly drain flow moderately well (NSE 0.57--0.64) and reproduced atrazine transport, macropore number and partitioning being key \parencite{Malone2014,Singh1996}; a modified version allowing redistribution and drainage during rainfall improved daily drainage at Iowa and Ontario sites to within 6--7\% cumulative \parencite{Xian2017}. \textsc{WEPP} enhancements improved uncalibrated drainage (daily NSE 0.27--0.47, rising to 0.59--0.73 calibrated) and reproduced storm and peak runoff around drainage installation \parencite{RevueltaAcosta2021,SAVABI1993}. \textsc{ADAPT} calibration reduced cumulative-drainage under-prediction from 14\% to 7\% in cold Minnesota \parencite{Mendez2004}, and \textsc{GPFARM-Range} simulated annual drainage acceptably (NSE $=0.76$) on drained pasture \parencite{Chen2016}. The \textsc{SIDRA} family, based on few parameters, has over three decades demonstrated predictive capacity for peak flows and midspan water tables and, when coupled (\textsc{SIDRA-RU}, \textsc{SIDRA-MAGE}), for catchment routing and solute transfer, with KGE $>0.81$ and a drain-flow-start median error of 5 days \parencite{Zimmer2023,Henine2022,LESAFFRE1988,Giraud1997}. Supporting these evaluations, automated logger--flowmeter systems measured drain flow to within 2\% of manual readings while resolving event dynamics \parencite{Kanwar1999}, and combined pressure-plus-discharge measurement diagnosed a tenfold decline in potential vertical flow in aged systems \parencite{Nishida2020,Gurklys2008}.

\subsection{Saturated versus Unsaturated Flow and Vertical Resistance}

A central mechanistic result is that treating drainage as purely saturated flow overstates discharge. A saturated-flow model over-predicted seepage and drainage time by roughly an order of magnitude relative to experiment, an analytical solution, and the SUTRA Richards-equation model, which agreed to within 20\% (seepage) and 40\% (drainage time); a vadose zone retaining substantial water above the phreatic surface persisted even at steady state, and performance depended strongly on capillary-fringe thickness, capillary rise, conductivity and drain elevation \parencite{Xin2016b}. Consistently, saturated--unsaturated drainage is governed by gravitational, pressure, hysteresis, capillary-barrier and air-pressure potentials, so saturated conductivity alone is an inadequate design basis; water-holding capacity, permeability and aeration jointly control efficiency, and preferential drainage can become discontinuous \parencite{Li2021}. Field diagnosis in Japanese rotational paddies showed that vertical, not horizontal, resistance limits the maximum drainage rate: potential vertical flow (2.1--5.7~mm~h$^{-1}$) was far below potential horizontal flow (6.1 to $>$79~mm~h$^{-1}$) \parencite{Nishida2020}. Accounting for vertical head loss improved paddy water-table prediction, differing from field profiles by 9--13\% versus 23--75\% for the conventional 2D solution \parencite{Talukolaee2022}. Neglecting unsaturated-zone buffering also underestimates response times: an analytical model retaining it matched a numerical model closely (steady heads within $0.6\pm2.0$~cm; real-case NSE $=0.80$ water level, 0.89 discharge), whereas ignoring the buffer distorted transients \parencite{VandeCraats2021}. Equivalent-medium and layered representations reproduce integrated outflow and mean water level well but can misrepresent internal dynamics, being least satisfactory where high unsaturated-zone conductivity permits fast vertical flow \parencite{Carlier2007}.

\subsection{Analytical Solutions for Drain Spacing and Water-Table Dynamics}

\subsubsection{Steady-State versus Unsteady-State Spacing}

Unsteady-state (transient) drain-spacing equations yield markedly narrower spacings than steady-state equations and are preferred where extreme short-duration water-table depths occur in high-storage systems \parencite{OOSTERBAAN1988}. In Iranian paddies, average computed spacing was 30~m for unsteady-state versus 55~m for steady-state equations, the unsteady solutions judged more suitable under variable rainfall \parencite{DarziNaftchally2014}. Steady Hooghoudt-based estimates matched observed spacings poorly in several settings, while among transient options the Glover--Dumm and van Schilfgaarde models were optimal for different rainfall durations \parencite{Moshayedi2020}. Design-criterion choice matters: for a 25~cm drop in 2 days a new equation estimated spacing closest to the actual 50~m (3.76\% deviation), whereas for longer periods the Glover equation was superior \parencite{Pali2014}. Nondimensional analyses show discharge following a common S-shaped curve against spacing ratio, with the median at a spacing-to-depth ratio of 2--5 \parencite{Nédélec2010}.

\subsubsection{Evapotranspiration, Recharge, and Sloping or Layered Aquifers}

Incorporating evapotranspiration and unsaturated-zone flow generally permits wider, more economical spacing. Analytical ET-aware solutions lower the water table faster and allow spacing increases of about 9--18\% (at $K=3$~m~day$^{-1}$, ET $=8$~mm~day$^{-1}$), applicable even to layered profiles \parencite{NIKAM1992,Singh2010,Upadhyaya2000}. Accounting for bare-soil evaporation widened calculated spacings by 23--74\% depending on texture (e.g. clay from 22.5 to 28~m) \parencite{AbdelFattah2005}. A spacing equation including horizontal unsaturated-zone flow gave up to 30\% wider spacing for coarse soils near an impervious layer, diminishing to $<$5\% when that layer lay 3~m below drain level, and agreed with HYDRUS-2D for non-submerged drains \parencite{Yousfi2014}. Bi-level solutions show the deep drain discharges more than the shallow one and that spacing is only slightly smaller than an equivalent conventional level system \parencite{Verma1998,Upadhyaya2000,Kankarej2023}. For sloping aquifers, water-table profiles become asymmetric with the divide shifting toward the lower drain, and hybrid finite-analytic solutions yield 12--17\% larger spacing than analytical ones \parencite{Singh2010,Upadhyaya2013,Hartani2001}. Dual-depth systems always remove water more effectively than single-depth systems (relative discharge 160--270\%), so retrofitting intermediate drains can raise effectiveness by 50\% or more \parencite{Kirkham1997}.

\subsubsection{Layered Soils, Ponded Drainage, and Benchmark Verification}

Analytical solutions for multi-layered and ponded-ditch systems have been verified against experiments, prior analytical work and MODFLOW, establishing that ponded drainage is strongly three-dimensional near the drains and highly sensitive to directional conductivities and anisotropy \parencite{Chakrabarti2020,Sarmah2015,SHARMA1991,Afruzi2014}. A plow-sole layer greatly restricts seepage and lengthens particle travel times, which can be reduced by staggered ponding and lower ditch levels \parencite{Chakrabarti2020,Sarmah2015}, and even a thin lower layer's conductivity significantly affects spacing \parencite{SHARMA1991}. Time-stepping linearisation of the Boussinesq equation outperformed earlier Werner and Baumann techniques \parencite{Lin2025}, and iteration schemes for non-homogeneous soils reproduced field water-table responses adequately \parencite{CHU1977}. On slopes up to about 0.15~m~m$^{-1}$, flat-land theory remains applicable for interior parallel drains, with deviations only at steeper slopes \parencite{FIPPS1989}; cross-model verifications show water-table shape-factor differences below 4\% for flat lands, with numerical instability for steep slopes \parencite{LORRE1994,Nédélec2010}. Critiques note, however, that a single linear-diffusion model cannot describe all stages of drawdown, since flow is radial near the tubes at early time and quasi-horizontal at late time \parencite{Kacimov2000}.

\subsubsection{Governing Equations and Single-Tile Flow}

The most physically complete descriptions solve the Richards equation, which fully represents saturated and unsaturated movement during drainage and subirrigation \parencite{TANG1977}. In sand-tank validation, predicted water-table drawdown and inflow/outflow rates and cumulative volumes agreed closely with observation, and, critically, approximate Boussinesq-based solutions agreed with the Richards and experimental results \emph{when a variable (nonconstant) drainable porosity was considered}; the Richards equation is more accurate, but corrected Boussinesq approximations remain useful for design \parencite{TANG1977}. Using an effective specific yield to mimic the vadose-zone effect, a new Boussinesq analytical solution removed large errors that arise when the unsaturated zone is neglected and gave good agreement with SUTRA and experiment for total drainage, though it predicted a quicker water-table decline and could not capture the dynamic drainable-porosity effect \parencite{Xin2016c}. Water-table recession following recharge begins only after vertical percolation ceases, with radial resistance appearing as a modified equivalent depth \parencite{VANDELEUR1959}, and classical transient-recharge Boussinesq solutions provide the theoretical basis for these dynamics \parencite{SINGH1980,Bouarfa2000}.

For flow to a single tile, analytical solutions closely agreed with finite-element simulations and, by replacing drain spacing with a variable \emph{region of influence}, became applicable to irregular or randomly spaced systems while eliminating the spacing parameter (RMS deviations 0.060--0.085~cm$^2$~day$^{-1}$) \parencite{Cooke2001}. Modified steady-state solutions of Ernst, Dagan and van Beers predicted transient midspan heads reasonably at 50--84~m spacings \parencite{SINGH1992}, and parameter-estimation methods agreed with Glover--Dumm to within 10\% for up to 15 days while providing high observed--predicted correlations (0.945--0.993) for the transmissivity-to-drainable-pore-space ratio \parencite{KUMAR1994}. A recurring obstacle is the drainable porosity: the Safaa equation avoids it via a new constant, giving spacing closer to actual and more economic than Glover's in the Nile Delta \parencite{FathAllah2025}, and an explicit minimum-cost ditch-spacing equation matched an optimisation program to within 7\% \parencite{Chahar2010}. Space-fractional Boussinesq formulations reduce to Glover--Dumm in homogeneous soil but track observed profiles more closely in heterogeneous soil, being most sensitive to the heterogeneity parameter, with elliptical forms fitting best \parencite{Mehdinejadiani2013,Mehdinejadiani2020}; a Bees-Algorithm inverse model gave repeatable conductivity estimates and the best simulation performance across datasets, improving on the Glover--Dumm criterion up to 34-fold \parencite{Mehdinejadiani2022}. Inverse frameworks using water-table-profile variation over time let all analytical models agree well with data, with variable effective porosity improving precision and the Topp--Moody model most accurate \parencite{Samani2007}; a Green-element solution resolved slope effects efficiently (2.5\% error at 30\% slope) \parencite{Onyejekwe1999}, and field determination yielded drainable pore space of 7.0--7.2\% and reaction factor 0.599~day$^{-1}$ consistent with reference values for silt loam \parencite{Bahceci2008}.

\subsection{Pipe-Scale Hydraulics and Structural Design}

The ease with which water crosses the pipe wall is characterised by the effective radius and entrance resistance. For Michigan pipes, 3--4-row regular-perforated pipes averaged 0.6~cm effective radius and 8-row reached 1.9~cm, while sand-slot pipes were lower; wrapping regular-perforated pipes in a knitted sock raised the effective radius dramatically to 5.7--6.0~cm and reduced entrance head loss to near zero \parencite{Ghane2022b}. Increasing the number of perforation rows was the most efficient way to raise effective radius, followed by perforation length, whereas width had minimal effect \parencite{Ghane2022b}. Sock-wrapped pipes allowed wider optimum spacing but higher initial cost than 8-row sand-slot pipes \parencite{Ghane2022b}; a companion study measured a sock-wrapped pipe's effective radius for the first time, validated the Dierickx theoretical method, and found the sock pipe functioned almost as a completely open pipe, with water entering only through bottom-half perforations under ponding \parencite{Ghane2022}. Numerical modelling calibrated to sand-tank experiments produced the first computation of entrance resistance for corrugated pipes with circular holes: because streamlines converge in the corrugation valley, resistance was higher (and delivery ratio lower) than in plain-wall pipes, and 51\% higher for circular holes than rectangular slots, making slots more hydraulically advantageous \parencite{Gaj2020}. Experimental validation of Dierickx's model confirmed that a screen envelope reduces the resistance factor and increases effective radius, that unenveloped pipes clogged almost instantly, and that smaller-diameter pipes with higher perforation density and a screen provided greater discharge capacity \parencite{Oyarce2017}. Bimodal (dual-mode) retention and conductivity functions predicted preferential flow to tile drains far better than unimodal van Genuchten--Mualem functions \parencite{Mohanty1997}, and a general discharge formula for improved drains under ponding was verified by experiment, field test and simulation \parencite{Tao2019}.

The structural counterpart shows that ductile yielding of non-perforated corrugated HDPE pipes is highly unlikely under field-traffic, soil and water-table conditions, with deflection within the 5\% limit; but circular-hole perforations substantially raise local wall stresses (5.6~mm holes exceeded yield strength shallower than 0.9~m in silt loam), whereas rectangular slots with width $\le 0.5w_v$ produce lower stress concentrations and lower failure risk---again favouring slots over holes \parencite{Gaj2021}.

Within-conduit flow with lateral inflow is a spatially varied flow (SVF) problem. An inverse-calibrated SVF model (Manning $n\approx0.02$) reproduced water-surface elevations, and ignoring local head loss under-estimated by up to 20\%; modified traditional relationships gave capacities two to four times the physically based SVF value \parencite{Bayat2011}. A generalized numerical solution for circular conduits established that Manning's equation alone is inadequate for drain capacity and that supercritical flow should be avoided \parencite{Graber2004}. Resistance depends on regime: reported Manning coefficients were 0.025/0.020 (GVF) and 0.035/0.030 (SVF) for two diameters, SVF increasing resistance 40--50\% over GVF \parencite{Kouchakzadeh2011}. Outlet studies identified ten flow patterns and non-negligible tail-water effects that can trigger free-surface-to-pressurised transitions \parencite{Nédélec2008}, and a 3D FLOW-3D model of an in-channel control module matched experiments to within 5\% while conserving mass \parencite{Abdurrasheed2019}. At watershed scale, a WASCoB routing toolbox reproduced flow and sediment through risers and berms, finding a 150~mm tile inadequate for design storms and recommending a 375~mm pipe \parencite{Gupta2019}. CFD of corrugated-pipe geometry (outside diameter dominant) matched sand-tank experiments to within 6\% \parencite{Wu2022}.

\subsection{Preferential Flow and Model Structure}

At field to hillslope scale, the representation of macropores and preferential flow is decisive for reproducing rapid drainage. Dual-permeability and dual-porosity models reproduced drainage hydrographs well (calibration $R^2=0.84$), the dual-permeability variant giving the best groundwater match ($r=0.98$), while a single-porosity model reached acceptable accuracy with far fewer parameters; all approaches failed to match low continuous dry-season discharge \parencite{Varvaris2018}. Dual-permeability formulations enable low matrix conductivities while still reproducing fast measured events through macropore flow, suiting complex topography and structured soils \parencite{Warsta2014,Salo2017}. In clay-dominated hillslopes, evapotranspiration and macropore flow dominate the water balance while Darcian matrix flow is negligible, though models can struggle to dry the profile in late summer \parencite{Koivusalo1999}. Adding macropore flow in MIKE SHE did not change total drainage but improved peak-flow simulation, with the drainage time constant (best value 5.6~h) most critical \parencite{Zhou2013}. Fully coupled surface--subsurface--tile models are recommended for field-scale dynamics, with discrete pipes most relevant for drain-flow dynamics and dual-permeability meshes reserved for detailed contaminant transport; representing the network as an equivalent porous medium or via seepage nodes can reproduce outflow almost exactly even when internal heads differ, offering a practical simplification for watershed integration \parencite{DeSchepper2015,Carlier2007,Boico2022}. Letting drainage occur simultaneously with rainfall and redistribution improves daily and peak-timing skill \parencite{Xian2017}.

\subsection{Tile Drainage Contributions to Streamflow}

\subsubsection{Annual and Event-Scale Partitioning}

Watershed-scale work establishes that tile drains route a large, measurable share of precipitation to streams. In an Ohio watershed, 21\% of precipitation (206~mm) was recovered through tile drainage, accounting for 41\% of annual and 47\% of monthly watershed discharge \parencite{King2014}. Equivalent-medium modelling of an Iowa watershed attributed 97~mm of tile drainage---13\% of precipitation and 39\% of streamflow \parencite{Thomas2016}. In the Boone River, end-member mixing (EMMA) and SWAT bracketed the tile contribution at roughly 46--54\% of discharge, SWAT estimating about 53.7\% (Boone) and 65.8\% (Lyons Creek) of annual yield \parencite{Schilling2019}. Snowmelt-dominated SWAT modelling found tile flow near 16\% of annual precipitation (about 37\% of yield), with tiling reducing surface runoff and increasing subsurface flow and yield \parencite{Rahman2014}. Field-scale export can be more tile-dominated: tile discharge accounted for $83\pm20$\% of Ohio water export \parencite{Williams2018}, and most discharge from an Indiana closed depression was delivered by tiles rather than surface flow \parencite{Williams2019}. Partitioning varies with hydrologic condition and area: in Iowa tile drainage averaged 30\% of streamflow during events but 61\% during non-event periods, event $Q_T/Q$ rose with drainage area while non-event $Q_T/Q$ fell, and annual $Q_T/Q$ rose log-linearly from 0.29 at 6~km$^2$ to 0.39 at 45~km$^2$ ($R^2=0.98$); the equivalent conductivity representing this scaled with drainage density from 164 to 530~m~d$^{-1}$ \parencite{Thomas2016,King2014}.

\subsubsection{Seasonal Dynamics}

Tile contributions peak when soils are wet and evapotranspiration is low. Ohio tile discharge was strongly seasonal (median monthly summer 0.9~mm versus winter 23.4~mm) \parencite{King2014}, and in the Boone River drainage dominated flow from March through June \parencite{Schilling2019}. \textsc{DRAINMOD} simulations indicated the largest baseflow effect from March to July---over 40\% more baseflow in drained watersheds and about 60\% more in April--June \parencite{Schilling2008}. Expanded drainage was projected to decrease winter flows and increase late-summer and fall flows \parencite{Rahman2014}. In cold regions tile flow tracked snowmelt with a seasonal absence from May to September \parencite{Kompanizare2024}, and water-table position relative to tile depth gates drainage---in Iowa a well-drained block's water table lay below tile depth 72\% of the time, steepening gradients during wet periods \parencite{Logsdon2020}.

\subsection{Modification of the Flow Regime}

\subsubsection{Low Flows and Baseflow}

A robust, repeated finding is that tile drainage elevates low flows and the baseflow fraction. Iowa modelling showed drained landscapes carried far higher baseflow (150~mm, 72\%, at 24~m spacing) than undrained (52~mm, 25\%), tile drainage primarily affecting the baseflow portion of the hydrograph \parencite{Schilling2008}. Ohio simulations found tiles greatly increased low flow while decreasing higher flows less, raised annual average daily flow about 5\%, and lowered the probability of the water table entering the root zone from about 50\% to about 5\% \parencite{Rathore2025}. Routing analysis reinforced that tiles increase low flows, decrease intermediate flows, and minimally affect the largest flows \parencite{Sloan2017}. In Latvia, drainage raised mean summer--autumn runoff by 132--186\% and winter low-water runoff by 53--286\%, with annual mean, minimum and maximum discharge up 46\%, 158\% and 26\% \parencite{Apsite2017}. Channelized Ohio headwaters showed far fewer zero-discharge days (8 versus 87.5~d~yr$^{-1}$ in natural streams), inverting part of the accepted headwater paradigm \parencite{King2009}.

\subsubsection{Peak Flows and Flood Response}

Effects on the largest peaks are consistently modest but scale- and condition-dependent for smaller ones. Routing modelling concentrated peakflow reductions at intermediate scales (100--1000~km$^2$), falling from 26\% at watershed to 11\% at basin outlet for one event, with negligible field-scale effect on the largest floods; placement mattered, draining headwaters producing the greatest reduction \parencite{Sloan2017}. Field-scale Iowa scenarios showed no significant difference in mean annual peak discharge for any of six soil types, indicating drainage chiefly alters pathway rather than the largest peaks \parencite{Sloan2016}. Ohio ATS simulations found tiles could enhance event peaks for sub-threshold or dry-condition events but reduce them for larger events, reducing streamflow in 23 of the 25 largest events \parencite{Rathore2025}. Under expanded tiling, smaller peaks became more frequent while larger peaks declined, steeper rising limbs implying reduced large-event peaks arose from reduced volumes rather than altered timing \parencite{Rahman2014}. Distributed modelling similarly found drainage decreases peak flow while increasing recession volume, reproducing observed behaviour with only 21\% of the area tiled \parencite{Fonley2021}.

\subsubsection{Flashiness, Drought and Long-Term Trends}

Flashiness responds to drainage in a soil- and slope-dependent way. A global meta-analysis of 464 watersheds found overall flashiness (Richards--Baker Index) similar regardless of drainage extent but diverging by setting: in flat watersheds ($<$2\% slope) medium/high drainage raised RBI, whereas in sloping watersheds low extent had the highest RBI, and drainage extent correlated positively with agricultural land use and clay content and negatively with slope \parencite{Zhou2025}. Field-scale Iowa simulations showed drainage decreased flashiness for low-conductivity soils and increased it for high-conductivity soils, homogenising response, with an optimal spacing near 15~m \parencite{Sloan2016}; Ohio modelling found overall flashiness decreased with tiles, minimised near 50~m spacing \parencite{Rathore2025}. Across 30 Nordic--Baltic catchments, flashiness related negatively to baseflow index and positively to drained share and intensity, and computing it from hourly rather than daily discharge increased it nearly threefold \parencite{Deelstra2014,Deelstra2014b}. At the regime scale, drainage significantly increased streamflow-drought duration and intensity while reducing frequency, each 10\% increase in tile-drained area lengthening drought duration by 0.03~d, raising intensity 12\%, and lowering frequency 0.10 events~yr$^{-1}$ \parencite{Adelsperger2024}. Indiana trend analysis found increasing low- and medium-flow trends with weakened precipitation linkage in heavily drained areas, suggesting a drainage influence \parencite{Kumar2009}, while a synthesis of seven Midwestern rivers concluded that inter-annual variability in streamflow, baseflow and N load was controlled primarily by current- and previous-year precipitation \parencite{Baeumler2020}.

\subsection{Scale Dependence, Homogenization, Travel Times, and Pathways}

Artificial drainage homogenises catchment response and reduces the scale dependence of hydrologic signatures. Non-tiled Iowa watersheds show decreasing area-normalized peak discharge and increasing baseflow ratio with scale, whereas tiled basins behave scale-invariantly; within-basin variability in recession coefficients declines as drained fraction rises, and the tiled/non-tiled distinction disappears above roughly 2200~km$^2$ \parencite{BolandBrien2014}. In Minnesota's Redwood basin, mean travel time fell 41\% (77 to 45~d), year-to-year variability was reduced, and travel time became independent of catchment size above about 200~km$^2$ \parencite{DaneshYazdi2016}; streamflow dynamics became more linear after 1976 despite little change in precipitation \parencite{FoufoulaGeorgiou2015}. Isotope and transit-time work in a UK catchment found short mean transit times ($<$1~yr, faster in artificial drains) and concluded that drainage homogenised distributed responses \parencite{DimitrovaPetrova2020}, and an Indiana watershed behaved as an effectively homogeneous unit describable by a single exponential travel-time distribution with a 2--3~d mean \parencite{Basu2010}. Even so, tile contributions remained detectable beyond 16{,}000~km$^2$ in the Boone River \parencite{Schilling2019}.

Modelling and tracers quantify how drainage shortens groundwater travel times. In Iowa, raising drainage density shortened mean travel time and raised the tile contribution to baseflow to an upper bound of 37\%, while increasing incision from 0.3 to 2.7~m lowered mean travel time near-linearly and raised tile contribution from 30\% to 54\% \parencite{Schilling2012}. In Bear Creek, mean groundwater travel time fell from 5.57 to 1.07~yr at high drainage intensity (versus 181.6~yr presettlement), short-circuiting flow to riparian buffers and cutting buffer residence time to 15.5--36.7~d, lowering nitrate-reduction effectiveness \parencite{Schilling2015}. Analytical solutions confirm that under tile drainage groundwater levels are lower and decline faster, total groundwater discharge to the river increases, and the tile contribution grows with density, though direct groundwater discharge during recession is reduced \parencite{Liang2015}; recession of heavily tiled watersheds is correspondingly more linear \parencite{Schilling2008}.

Runoff generation is governed by pathway-activation thresholds. In an Indiana closed depression, rainfall alone related linearly to cumulative flow with no threshold, but a clear threshold emerged when antecedent soil-moisture deficit and rainfall were summed, varying by pathway and landscape position; tile flow generally began earlier, peaked earlier and lasted longer than surface flow \parencite{Williams2019}. Isotopic work there showed tile celerity was controlled by antecedent wetness, with preferential flow more important when dry and stored pre-event water dominating when wetter, and groundwater the primary tile source (mean transit time to tile 245~d) \parencite{Williams2021}. In near-level Canadian vertisols, overland flow generally activated before tile flow, tile activation relating to maximum rainfall intensity and antecedent precipitation, with direct preferential flow into tiles uncommon \parencite{Kokulan2019}. Across 795 US--Canadian events, 43\% produced both surface and tile flow and 57\% tile flow only, surface flow contributing a median 12\% among dual-pathway events, with event-response types partitioning by texture and season \parencite{Macrae2019}. In Kansas tile-outlet terraces, baseflow contributed 38--51\% of discharge, antecedent conditions and storm size being stronger drivers than tile-inlet density \parencite{Stops2022}.

Tracer and geochemical studies partition tile flow into fast preferential and slower matrix components whose balance shifts with season and wetness. In a French field, dry-condition rainwater transferred directly through macropores whereas saturated conditions were dominated by pre-event soil water, with surface and tile drainage transferring 44\% and 20\% of annual effective rainfall and 65\% of infiltrating water not captured by tiles \parencite{Gaillot2023}; catchment monitoring identified slow saturated drainage and fast preferential flow, the fast peak contributing about 15\% of water and sediment fluxes \parencite{Grangeon2021}. Ohio master-recession analysis distinguished two reservoirs and apportioned annual tile discharge into 9\% quick-old, 39\% quick-new, 42\% slow-old and 10\% slow-new, showing new water routes through both reservoirs \parencite{Nazari2021}. Preferential flow accounted for a higher fraction in clay than loam and correlated with phosphorus loads \parencite{Pluer2020}, and Swedish modelling attributed autumn events to preferential flow and spring events to matrix flow \parencite{Salazar2011}. Flow and export are highly concentrated: in Ohio, surface runoff occurred $41\pm42$ and tile drainage $240\pm69$~d~yr$^{-1}$, 80\% of tile discharge in 21\% of the year; a 30\% cut in annual tile NO$_3$-N load would require treating tile flow 83\% of the time if targeting low flows but only 3\% if targeting high flows \parencite{Williams2018}. Managed-drainage monitoring reinforced this---287~mm (40.3\% of precipitation) of discharge, 94\% when the valve was open \parencite{Miller2022}.

\subsection{Watershed- and Regional-Scale Representation}

How drains and subsurface heterogeneity are represented strongly controls simulated catchment behaviour. Two grid conceptualisations of a Danish catchment produced drain contributions of 13\% versus 31\% despite both matching total streamflow within 0.1\% \parencite{Karan2021}. In seepage-node conceptual models, soil heterogeneity affected drainage-area simulations more than tile conceptualisation, while the reverse held at subcatchment scale; seepage-node models were about 35\% faster and better suited to automated calibration and unknown drain locations \parencite{Boico2022}. Distributed HydroGeoSphere models reproduced flashy heads and hydrographs better than MIKE SHE \parencite{DeSchepper2017}, and machine learning established that streamflow signatures (especially high-flow duration) predict simulated drain fraction across 73 Danish sub-catchments (test MAE $=0.126$) \parencite{Schneider2022}. Novel strategies embedding drainage in fully distributed integrated models reproduced streamflow without site-specific calibration, outperforming a calibrated SWAT model \parencite{Rathore2024}, while 3D CATHY simulations showed drainage raising cumulative outlet flow 3--22\%, lowering peaks and raising minima \parencite{Muma2016}. Water balances quantify drain contributions: subsurface drainage was the dominant source (about 68\% of total drainage) in a Chinese irrigation district \parencite{Zheng2025b} and a major precipitation fraction across Nordic catchments \parencite{Nousiainen2015,Turunen2013}. Equivalent-medium representations reproduced explicit tile simulations within 2\% of drained-volume targets \parencite{Thomas2016}, physically based cold-region modelling captured tile timing and the seasonal absence of summer flow \parencite{Kompanizare2024}, a modified ODE hillslope formulation improved recession and flood-threshold timing at drained gauges \parencite{Fonley2021}, and SWAT was tested against daily tile flow (NSE $=0.5$, PBIAS $=-1.4$\%) \parencite{Rahman2014}.

Distributed groundwater models, principally MODFLOW, are established tools for regional assessment, salinity-hazard mapping and optimisation. Calibrated models matched observed heads closely (e.g. 92\% correlation) and were used to design interceptor systems, quantify drawdown (average 0.53~m over three years) and predict discharge \parencite{Ghorbani2016,Mirlas2012,Mirlas2009,Mirlas2013,Mirlas2021,AbdElaziz2025}. Sensitivity analyses showed these models relatively insensitive to specific storage ($\pm$15\%) and aquifer conductivity ($\pm$50\%) but responsive to aquitard conductivity, and revealed existing drains sometimes intercepting only 9--30\% of groundwater flow \parencite{Mirlas2012,Mirlas2009,Mirlas2013}. Simulation--optimisation achieved equal or better water-table control with about 20\% less pumping \parencite{Rana2008}, a 90$^{\circ}$ alignment to groundwater flow was identified as optimal \parencite{AbdElaziz2025}, and coastal applications quantified sea-level rise raising drain discharge by 12.5--43.4\% \parencite{AbdElaty2019}. New MODFLOW subsurface-drainage (SDR) packages based on the Hooghoudt equation markedly outperform the conventional DRN package, reproducing steady-state and experimental drainage (RRMSE $<$5\%) by calibrating only conductivity, whereas fixed DRN conductance caused large errors and simulated 16.9--55.7 times more drainage than DRN \parencite{Yang2022b,Han2025}; an entrance-resistance modification corrected over-predictions for synthetic-envelope pipes (NRMSE $<$3.18\%; NSE, KGE $>$0.9) \parencite{Qian2025b}. Such frameworks quantified regional salt-discharge increases of 38--54\% and vadose-zone salt reductions exceeding 90\%, a critical drainage-to-irrigation ratio near 10--11\%, and a 1~m appropriate grid resolution \parencite{Yang2022b,Han2025}. Regional nitrate mapping showed high-drainage areas tend to have moderate-to-low groundwater reduction because drainage bypasses reducing zones \parencite{Hojberg2017,Hansen2019}.

\subsection{Controlled and Managed Drainage}

Controlled and managed drainage strongly modulate outflow volume and timing. In Sweden, controlled drainage reduced total outflow by 73--95\% relative to conventional drainage, concentrated more than 60\% of daily outflow in the 0--0.5~mm class, and produced lower retention constants and near-zero final storage \parencite{Wesström2003}. Raising the drainage outlet slowed water-table recession by 29\% and 62\% in two site pairs, confirming the effect was not seasonal \parencite{Saadat2017}. Agricultural drainage-water management in Ohio was most effective during low-intensity events and dry antecedent conditions, the valve open on average 18.6\% of each three-day event yet passing 96\% of discharge, though average groundwater and soil-moisture indicators did not differ significantly from free drainage \parencite{Miller2022}. Netherlands subsurface water-infiltration systems significantly influenced water-table depth, reducing fluctuations at four of five sites \parencite{vandeCraats2025}. Cross-drain spacing tunes recession, being most effective at 20~m (reducing recession 13--32\%) and negligible by 40~m \parencite{Vaiciukynas2015}, and trenchless installation held groundwater 0.14--0.25~m higher than trench-installed sections in 60--90\% of observations \parencite{Salo2019}. Method choice matters for evaluating such experiments: paired-watershed, BACIP and distribution-based approaches can yield differing conclusions, a beta distribution deemed most appropriate in one Minnesota dataset \parencite{Pagliari2021}. Broader analyses established controlled drainage's water-quality benefits---raising groundwater tables 19--35~cm and reducing annual outflow 35--52~mm and nitrate losses 15--22~kg~ha$^{-1}$ \parencite{Kesicka2023}---while lateral seepage explained 96\% of the drainage reduction, so ignoring seepage severely over-predicts outflow \parencite{Youssef2021,Saadat2018,Österholm2015}. Bibliometric analysis of 462 articles confirms controlled drainage's emergence as a leading multidisciplinary water-quality practice \parencite{Kesicka2023b}, and emerging model-predictive-control and drainage-water-recycling tools extend capabilities toward real-time water-table management \parencite{Ikonen2023,Reinhart2020,VandeCraats2021}.

\subsection{Drain Hydraulics, Sediment, Solute, and Specialised Settings}

At the pipe scale, network pressurization introduces storage and delayed secondary peaks. In France, pressurization of the collector network during intense rainfall limited and sometimes reversed discharge, temporarily storing recharge as the water table rose and generating a four-stage sequence that increased the lag to the catchment outlet to 10--25~h \parencite{Henine2010}; a coupled groundwater--pipe model reproduced this, with secondary peaks 11--30~h after rainfall and storage of 11--51~m$^3$ (38--46\% of stored recharge) \parencite{Henine2014}. Drain flow is governed by the system parameter $\sigma=K/(f^2L^2)$, maximum flow rising from 0.025 to 0.56~mm~h$^{-1}$ as $\sigma$ increased \parencite{Bouarfa2000}. Combined tile/mole systems modelled with SEEP/W reproduced discharge well ($R^2=0.811$), a tile-only system being least effective (63\% less drain discharge than the best combined system) \parencite{Tuohy2018}, and drain-flow spatial patterns were driven by topography and geology \parencite{Mahmood2023}. Component-scale studies show tile risers and blind inlets altered flow duration without necessarily changing discharge frequency, blind-inlet infiltration declining to an 8--10-year service life \parencite{Williams2020,Williams2019}; blind or gravel inlets substantially reduced sediment and phosphorus losses (TSS 14.4 versus 40.4~kg~ha$^{-1}$~event$^{-1}$) with a flow-capacity/sediment-control trade-off among designs \parencite{Feyereisen2015b,Li2017b}. In recycled-concrete media, permissible flow fell up to 73\% at 4\% fines \parencite{McCulloch2017}, sheet-pipe drainage drained about 45\% more water with 55\% lower percolation \parencite{Tamura2021,Arif2021}, and two distribution pipes could raise 10-day cumulative infiltration in saturated buffers by up to 133\% \parencite{Jaynes2019b}. Deep-percolation recharge has long been recognised as part of the drainage water balance \parencite{BOOKMAN1959}.

Because streamflow pathways carry material, hydrologic partitioning shapes sediment and solute export. In dual-drainage French fields, mean suspended-sediment yield was 0.79~t~ha$^{-1}$~yr$^{-1}$, 89\% from surface drainage, tile drainage contributing 46\% and 21\% of runoff but only 9\% and 11\% of sediment \parencite{Gaillot2023b,Grangeon2021}. A sediment-connectivity index captured surface sources well but not subsurface-tile-controlled vertical flow \parencite{Calsamiglia2018}. Solute transport tended toward chemostatic behaviour dominated by antecedent conditions and storm size \parencite{Stops2022}, tile drainage delivered elevated antibiotic-resistance genes correlated to flow rate \parencite{Rieke2018}, and converting surface runoff to subsurface drainage did not necessarily improve water quality, a drainage-type raingarden retaining only 36\% of flow and sometimes exporting more nutrients than it received \parencite{Xu2022}. Beyond humid croplands, subsurface drains lowered surface electrical conductivity and suppressed resalination in saline Chinese soils \parencite{Yan2013,Chang2012}, passive drains stabilised Ethiopian gully banks (cutting bank retreat 52--84\% and soil loss about 70\%) \parencite{Zegeye2021}, and under-phreatic permeability increased downstream in a Burkina Faso valley bottom, so no single infiltration process described the toposequence \parencite{Keita2014}. Coupled SC-EMMA and recession frameworks improved dissolved-phosphorus prediction \parencite{Nazari2021}, and pipe-flow modelling feeds edge-of-field practice design: balancing nitrate removal against travel time gave optimal saturated-buffer widths smaller than current widths, implying removal could be 50--120\% greater at four of six sites, with 57\% uncertainty dominated by the removal coefficient \parencite{McEachran2020}.

\subsection{Data-Driven, Remote-Sensing, and Geostatistical Methods}

Data-driven methods frequently match or outperform process models for the specific outputs they are trained on while requiring less data and computation. Artificial neural networks predicted real-time water-table control accurately (minimum RMSE 60.6--73.8~mm), most best networks being small \parencite{Yang2000}, and predicted saturated conductivity more accurately than pedotransfer functions ($R^2$ up to 0.96), with ANN-based designs closer to measured-data designs and trainable on about 50 measurements \parencite{Ellafi2021}. Random forests outperformed multiple linear regression for agronomic prediction (water-table depth $R^2=0.95$ versus 0.28), capturing nonlinear effects and locating an optimal 7~m spacing \parencite{Djiemon2019}. An Attention U-Net surrogate with adaptive weighted loss cut $L_2$ error about 43\% and reached test $R^2$ of 0.997--1 for head, drainage volume and salt discharge while drastically reducing computation \parencite{Lei2026}. For salinity, a back-propagation network predicted drainage-effluent salinity better than SALTMOD, but SALTMOD remained superior for root-zone salt balance, the ANN being a faster ``gray box'' on minimal input \parencite{Sarangi2006}. Clustering and discriminant analysis delineated management zones, identifying elevation and subsurface drainage as significant discriminators of NO$_3$-N leaching zones (86\% accuracy) \parencite{Bakhsh2005b,Bakhsh2007b}.

Remote sensing combined with deep learning now maps tile-drained extent at regional to national scales with accuracies commonly exceeding 85--96\% and totals within a few percent of census estimates. AgTile-US reported a national tile area of 22.34~Mha ($-0.35$\% versus USDA) with 86.03\% overall accuracy \parencite{Valayamkunnath2020}, and the SEETileDrain random-forest product achieved 0.96 testing accuracy and F1 of 0.90, estimating 185{,}549~km$^2$ (about 9.4\% below USDA-NASS) \parencite{Wan2024}; a Danish selective ensemble predicted drained area over 49.4\% of the region ($\kappa=0.526$) \parencite{Moller2018}. Convolutional segmentation networks reliably delineate visible drain lines: a U-Net exceeded 92\% accuracy including 200~km beyond training data, drains most visible in spring under high radiation \parencite{Redoloza2023}, and Ohio and ensemble segmentation models outperformed Sobel/Canny edge detectors and the Hough transform, with resolution-dependent performance peaking near 400~m/256~pixels \parencite{Song2021,Woo2023}. Earlier decision-tree classification estimated 50\% drained area with 43--44~m sub-watershed spacing \parencite{Naz2009}. Ancillary methods infer drainage-relevant surface conditions: SAR discriminated drained from undrained corn by contrast \parencite{Meunier2001}, Landsat SWIR estimated soil-water content ($r=0.94$) \parencite{Rijal2013}, SAVI estimated cover-crop biomass at 86--107\% of ground values \parencite{Xu2018}, and high-resolution thermal imagery can detect drain lines, though subsurface drainage remained comparatively under-studied \parencite{Khanal2020,Zhang2010}.

Kriging-family methods are consistently the most reliable interpolators for the spatially variable soil properties governing drain spacing. In the Nile Delta, hydraulic-conductivity variation was isotropic with a 5.6~km range, kriging cross-validation giving low errors, and kriged conductivities recommended as design values \parencite{GALLICHAND1991}; a 600~m grid retained nearly all information and a 900~m preliminary survey fixed the optimal density \parencite{GALLICHAND1992}. For clay content no method showed important bias, cokriging marginally best \parencite{GALLICHAND1993}. Accounting for variability has quantified value: measured water-table depth underestimated drainage need by about 40\% and measured salinity overestimated by 59\%, while kriged conductivity widened spacing about 5\% and saved 17{,}155~m of pipe while reducing sample size three- to sevenfold \parencite{Moustafa1998b}; neglecting $K_{sat}$ variability could bias design spacing by $-27$\% to $+3$\% \parencite{Moustafa2000}. The covariance variogram proved superior to the semi-variogram for clay soils, and subsurface drainage itself altered spatial variability, the correlation range increasing about 29\% four years after installation \parencite{Moustafa1998}. Ordinary point kriging suited $K_{sat}$ and soluble salts in Turkey (horizontal correlation 710--1130~m) \parencite{Akis2015,Sepaskhah2004}.

\subsection{Sensitivity, Calibration, and Uncertainty}

Formal analyses converge on saturated hydraulic conductivity and evapotranspiration estimation as the dominant sources of sensitivity and uncertainty. \textsc{DRAINMOD} was most sensitive to vertical conductivity of the restrictive layer and lateral conductivity of the deepest soil layer, parameter interactions appearing only under likelihood weighting \parencite{Wang2006}. Conductivity estimation method produced ``tangible'' differences in simulated water-table fluctuations \parencite{NabiSichani2012}, and the van Genuchten shape parameter $n$, lateral conductivity and depth to the impermeable layer ranked as the leading sensitive parameters in HYDRUS inverse modelling, where monitoring discharge together with water-table depth was sufficient (NSE up to 0.83--0.95) \parencite{Sedaghatdoost2018}. GLUE/extended-FAST frameworks showed annual predictions within confidence bounds while exposing interaction effects invisible to one-at-a-time analysis \parencite{Wang2006}. Multi-objective sensitivity analysis of 17 parameters found daily drainage flow and relative corn yield most sensitive to drain spacing and depth \parencite{Timalsina2023}, the soil-water control volume was most influential for drainage-water-recycling design \parencite{Reinhart2020}, and parameter uncertainty propagated materially into predicted root-zone salinity and crop yield, motivating probabilistic design \parencite{Alzraiee2013,Alzraiee2013b}.

The construction and optimisation of objective functions strongly affects identifiability. A genetic algorithm generally outperformed PEST, combining discharge, water-table depth and salinity being best, though discharge with water-table depth alone already gave excellent results \parencite{Sedaghatdoost2019}. Variational data assimilation---the first for SIDRA-type models---identified $K_{sat}$ as most influential, selecting only 4 of 11 parameters, and cut the five-year cumulative-volume discrepancy to $<$1\% versus 10\% for a prior method \parencite{Chelil2022b}. Regional single-objective calibration in California produced identifiable parameters and predicted water-table variation well (RMSE 2.76~m), multi-objective calibration revealing trade-offs between water-table and pumping predictions, with drainage predictions more sensitive yet largely independent \parencite{Schoups2005}. The review literature concludes that no universal model exists, that automatic calibration is preferable because manual calibration neither guarantees global optima nor supports uncertainty estimation, and that integrated multi-model use is often more practical \parencite{Singh2019b}.

\subsection{Evaluation Metrics and Measurement Protocols}

The choice of goodness-of-fit metric materially changes apparent performance and even model ranking, so metric selection is itself a substantive modelling decision. Across 14 Illinois datasets the Kling--Gupta efficiency consistently outperformed the Nash--Sutcliffe efficiency, with the same behaviours recurring across every dataset, while the heteroscedastic maximum-likelihood estimator ranked between them but depended strongly on the error structure, producing a ranking reversal for 4 of 14 datasets under one scenario \parencite{Bruhn2025}. In a paired storm-runoff example, simulations grossly over-predicted all 10 events, yet this was not reflected by $R^2$ or standard error; a regression-based test (intercept $\ne0$, slope $\ne1$) instead deemed the model unacceptable, and RMSE and mean bias error better reflected the discrepancies \parencite{STONE1995,SAVABI1995}. Curve-number estimation for 12 Ohio fields likewise showed correlation and NSE varying widely by method, the complacent--violent method giving the best fit and largest-event prediction \parencite{Wessel2022}.

Performance evaluation extends to the field protocols generating the validation data. \textsc{DRAINMOD} achieved good monthly agreement across five soil datasets, with ROSETTA-estimated $K_s$ driving it as accurately as measured values in coarse-textured soils (drainage errors generally $<$2\%) \parencite{Salazar2008}. \textsc{SIDRA-RU} simulated discharge accurately, especially on silty soils (calibration KGE$'$ 0.44--0.83), a split-sample test confirming national-scale robustness \parencite{Jeantet2021}, and HYDRUS-2D reproduced drainage nitrate very well (efficiency 0.6--1.0), estimating the grey water footprint within about 12\% of field values versus up to 3960\% error for a tier-1 method \parencite{Naderi2025}. For phosphorus-load monitoring, all sampling strategies underestimated load, but flow-proportional designs stayed within $-0.2$ to $-7.2$\% while time-proportional designs worsened to $-51$\%, giving more accurate estimates with fewer samples \parencite{Dialameh2022}. A recording ultrasonic flow meter achieved very high calibration quality ($R^2$ up to 0.9993) with $Q(h)$ uncertainties of 0.1--0.4\% \parencite{Enamorado2007}, while a continuously deployed bromide sensor under- and over-reported concentrations, post-hoc bottle measurement outperforming real-time deployment \parencite{deOliveira2023b}. Non-intrusive geophysics mapped as-built golf-green tile, and both large-ring and small-core methods were valuable for relative comparisons, infiltration being 5--65 times larger in trench backfill than midway between trenches in Swedish clays \parencite{Freeland2014,Messing2006,BUCKLAND1986}.

Accurate discharge measurement underpins model validation and conservation assessment. For a 45$^\circ$ V-notch weir in AgriDrain structures, the fitted $Q=0.011H^{2.28}$ ($R^2=0.992$) held with standard deviations of 0.1--2.9\% of the mean, valid to 16.5~cm depth \parencite{Christianson2019}; a metal-edge V-notch gave $Q=0.749H^{2.25}$ ($R^2=0.999$), higher than the AgriDrain equation \parencite{BinShokrana2021}, and a 22.5$^\circ$ weir's 45$^\circ$ chamfer gave the strongest fit ($R^2=0.9981$) \parencite{Napierala2023}. Because V-notch equations fail once flow overtops, a compound-weir equation estimated overflow within MAPE $<$5\% whereas the combined V-notch-plus-rectangular method under-estimated by over 15\%, materially biasing nutrient-load estimates diverted to a saturated buffer \parencite{Katuwal2024}. Adjustable damming regulators were calibrated for different flow ranges \parencite{Popek2021}, and a low-cost vision-based method detected outlet flow within 12\% of ground truth, releasing an open-source application \parencite{Young2023}.

Aggregative indicators and empirical regressions provide a complementary, policy-facing layer. A global compilation of 60 hydraulic-retention-time--nitrate-removal regressions showed strong within-study linearity but wide between-study variability, slope correlating positively with temperature and negatively with inflow N, and yielded generalised median models (NRE $\approx2.9\times$HRT$+0.37$ overall; $3.2\times$HRT$+0.18$ for a tile proxy) \parencite{Christianson2025c}. A soil quality index was unaffected by tillage or drainage but correlated with corn yield ($R=0.62$), $K_{sat}$ the key physical indicator \parencite{Nakajima2015}, while phosphorus-loss indicators and the Phosphorus Index are framed as source $\times$ transport aggregations expected to be directionally correct against water-quality data \parencite{Reid2018,Zhang2018}, and a riparian-quality index correlated with an independent measure, width alone correlating most strongly \parencite{Terrado2015}.

\subsection{Design Applications and Field Outcomes}

The validated models have been applied extensively to quantify design trade-offs. Narrower spacing and shallower drains improve water-table control and aeration: SEW30 fell sharply from 3461 to 450~cm~day across configurations as drains were made shallower and closer in Iranian paddies, a 0.65~m depth with 30~m spacing providing desired wet-season control \parencite{DarziNaftchali2013,DarziNaftchally2014}. Design analyses balance aeration against drainage volume, with SEW30 rising and drainage volume falling with wider spacing and mild-to-severe aeration risk at 75~m but not 25~m \parencite{Verma2014,Hirekhan2007}. Optimum spacings for trafficability on Illinois claypan soils were 10--13~m \parencite{MOSTAGHIMI1989}, and in South Africa target water tables of 1.0--1.5~m required 25--40~m spacing in clay versus 55--70~m in clay loam, confirming the closer-spacing-for-lower-$K_{sat}$ rule \parencite{Malota2015}. Shallow drainage cut Iowa subsurface drainage by 8.4\% \parencite{Zhou2013,Mendez2004}. Climate-change simulations consistently project higher future subsurface drainage (about 14--80\% depending on model and precipitation) concentrated in winter, with implications for nitrate loss, while recommended designs still maintained relative yields above 80\% \parencite{Singh2009,Chen2016,Feng2018}. Supporting field studies established the reliability of the measurement and property-estimation methods underpinning this modelling: drain-derived conductivity and drainable porosity agreed with independent methods \parencite{BUCKLAND1986}, Rosetta and RETC reproduced measured conductivity and retention accurately \parencite{Malota2022,Malota2024}, cover crops reduce drainage by a robust $\approx30$~mm \parencite{Meyer2019}, and long-term drainage intensity measurably alters soil physical properties, carbon accumulation and workability, with $\le12$~m spacing protecting subsoil from compaction \parencite{Omonode2025,Mante2022,Mante2018,Frison2009,Roy2020}. Together these results demonstrate that the hydrologic and streamflow signature of tile drainage---accurate monthly and annual water balances, elevated low flows and baseflow, damped intermediate flows, modest large-flood effects, shortened and homogenised travel times, scale-dependent flashiness and peak response, and quantifiable design trade-offs---is both mechanistically explicable and reproducible across observational and modelling lines of evidence \parencite{Ahmadi1995,HASEGAWA1994,GALLICHAND1992b,MCCARTHY1992,PalaciosVélez2004,Khan2007,Prathapar1996,Ghorbani2015,Skaggs2017,Singh2006,Singh1999,Gaj2022}.